% ---------------------------------------------------------------------------
% assignments.tex -- description of course assignments for ECON 7102.
%
% Hand-maintained. Unlike schedule_table.tex and reading_list.tex, this file is
% not generated from data/schedule.yml. Edit it directly.
%
% \input by econ7102_fall2026_public.tex.
% ---------------------------------------------------------------------------

\subsection*{Attendance and participation}
Attendance is expected and will be necessary to succeed in the course. Participation in class discussions is strongly encouraged. Ask questions. If you do not understand something in class, you are almost certainly not the only one. If you never volunteer an incorrect answer, you aren't taking enough risks.

\noindent Attendance is not graded directly. Poor performance in this dimension may affect the final grade.

\subsection*{``Weekly'' reading responses (20 points)}
Students write a short response paper, roughly 250--300 words or two paragraphs, on designated readings. These readings are marked \textbf{[response due]} in the reading list. Expect eight to ten over the semester.

\noindent The response needs to go beyond summarizing the reading. Make connections to other topics in the course, to stories in the news, or to ongoing policy discussions. Reflect on how the reading updated your thinking. If you disagreed with a conclusion or a fundamental assumption, expand on that.

\noindent Focus on three questions:
\begin{enumerate}
	\item What was one thing you enjoyed about the article?
	\item What was one thing you did not understand, or would like to discuss in class?
	\item What extensions or applications of the article would be interesting?
\end{enumerate}

\noindent Submit via Canvas by 11:59 PM ET the day before the class in which the reading is assigned. That timing lets me read the responses before class. Late responses receive half credit.

\subsection*{Research paper proposal (35 points)}
One goal of this course is to provide the tools necessary to conduct your own research in environmental, energy, and resource economics. Students develop a research proposal that introduces a novel and answerable research question, reviews related literature, discusses potential contributions, and begins to explore data sources and empirical methods.

\noindent There is no explicit page requirement, though 10--15 pages with normal formatting is expected. The proposal should look and read like an extended introduction to an economics journal article: title, abstract, main body, background where relevant, conceptual framework, and references. Because this course focuses on economic theory applied to environmental problems, bring that conceptual thinking to your proposal even if the goal is an empirical paper.

\noindent Research takes time, with many fits and starts, so there are milestones through the semester. Start thinking about ideas as soon as possible and discuss them with your classmates and with me. Early feedback prevents rabbit holes. Topics may need revision if they are not aligned with the course or do not represent feasible projects. Ideally, proposals become third-year papers or dissertation chapters.

\begin{center}
\begin{tabular}{llc}
\toprule
\textbf{Milestone} & \textbf{Due} & \textbf{Points} \\
\midrule
Project idea and research sketch		& October 9		& 5 \\
Expanded outline and literature review	& November 6	& 5 \\
Presentation							& Last four classes & 5 \\
Final proposal							& December 11	& 20 \\
\bottomrule
\end{tabular}
\end{center}

\noindent Students cannot use this proposal in conjunction with a writing assignment for another course. Proposals can sit at the intersection with other fields, for example environment and health, environment and development, or environment and trade.

\subsection*{Midterm (25 points)}
In class, Class 15.

\subsection*{Lead in-class discussion (10 points)}
Each student leads or co-leads discussion of one paper in the topics section of the course.

\subsection*{Referee report on a job-market paper (10 points)}
Students find a job-market paper by a current Ph.D. candidate working on an environmental, energy, or resource economics topic, and write a referee report. The goal is to see the end product of the Ph.D. on a topic of interest and to practice giving critical but constructive feedback.

\noindent I will provide guidance during the semester on where to search for job-market papers and on the structure of a referee report.
